\documentclass[11pt]{article}

\usepackage[utf8]{inputenc}
\usepackage[T1]{fontenc}
\usepackage{lmodern}
\usepackage{geometry}
\usepackage{amsmath,amssymb,amsfonts,amsthm,mathtools}
\usepackage{bm}
\usepackage{enumitem}
\usepackage{microtype}
\usepackage{hyperref}
\usepackage{titlesec}
\usepackage{booktabs}
\usepackage{array}
\usepackage{setspace}
\usepackage{mathrsfs}
\usepackage{physics}

\geometry{margin=1in}
\hypersetup{colorlinks=true, linkcolor=black, urlcolor=blue, citecolor=black}
\setlist[enumerate]{topsep=0.4em,itemsep=0.35em}
\setlist[itemize]{topsep=0.3em,itemsep=0.25em}
\titleformat{\section}{\large\bfseries}{\thesection.}{0.5em}{}
\titleformat{\subsection}{\normalsize\bfseries}{\thesubsection.}{0.5em}{}
\setstretch{1.06}

\newcommand{\Aether}{\AE{}ther}
\newcommand{\Aflow}{\AE{}ther-flow}
\newcommand{\TheoryName}{\Aflow{} Interpretation of Relativity}
\newcommand{\TheoryProgram}{\Aether{} / \Aflow{} framework}
\newcommand{\CompletedMetric}{g^{\sharp}}
\newcommand{\RootNucleus}{\mathfrak N^{\mathrm{rt}}}
\newcommand{\SeedPackage}{\mathfrak S^{\mathrm{seed}}}

\newtheorem{definition}{Definition}
\newtheorem{proposition}{Proposition}
\newtheorem{remark}{Remark}
\newtheorem{corollary}{Corollary}
\newtheorem{theorem}{Theorem}

\title{\textbf{The \TheoryName{}}\\[0.4em]
\large Primitive-Reservoir Observer-Reduction\\
\large Nucleus-Generating Substrate Seed Dynamics\\
\large Next Benchmark Criterion}
\author{}
\date{}

\begin{document}

\maketitle

\begin{abstract}
The active primitive-reservoir observer-reduction line has now moved one honest step further below the current root-nucleus frontier. The previous collapse theorem already showed that the larger benchmark-facing root core carries no extra benchmark-relevant freedom beyond the minimal root exact-shell transport nucleus \(\RootNucleus\). The new seed-dynamics theorem then removed the next unexplained stopping point available there: the nucleus itself is now the exact-shell transport image of a more primitive nucleus-generating substrate seed dynamics package \(\SeedPackage\).

This note records the routing consequence of that new theorem. Another same-output restatement of the minimal root exact-shell transport nucleus, another same-output root / foundation / ground relay, or a return to stationary-reduction, stationary-Hessian, spectral, or selector layers no longer counts as the right next benchmark-facing manuscript. The next honest benchmark-facing move must therefore either derive or justify the seed dynamics \(\SeedPackage\) from still more primitive substrate structure in a way that changes benchmark-facing status, or prove a genuinely smaller theorem-level collapse strictly inside the observer-relevant seed package itself.

The benchmark boundary remains fixed throughout. The one operative metric is still exactly \(\CompletedMetric\), the ordinary matter route is unchanged, there is no second operative metric, the low-energy matter law is not enlarged, and no stronger promotion language is licensed. The positive result of this note is therefore routing discipline at the new seed-level frontier.
\end{abstract}

\section{Introduction}

The active derivational program of \TheoryName{} is no longer at a stage where the minimal root exact-shell transport nucleus should be treated as the default unexplained stopping point. The previous root-nucleus collapse theorem had already removed the larger root-core freedom. The new seed-dynamics theorem then spent the next honest move available there by showing that the current nucleus itself is generated by a deeper nucleus-generating substrate seed dynamics package.

The present note records the routing consequence of that theorem. It does not reopen the seed mathematics, and it does not claim a completed first-principles substrate derivation of gravity. It records only which kind of next manuscript would now count as an honest benchmark-facing gain below the new seed frontier.

\paragraph{Framework and claim boundary.}
\TheoryProgram{} denotes the broader \Aether{} / \Aflow{} framework built on the claims that the \Aether{} is the underlying four-dimensional substrate of reality, the \Aflow{} is its intrinsic ordered motion, observed three-dimensional space is the local experiential slice of that deeper substrate, S-time is the experienced order of change arising from matter, light, and the \Aflow{}, and observed expansion is the three-dimensional appearance of deeper four-dimensional ordered motion. Within that framework, \TheoryName{} denotes the exact-closure theory in which the effective gravitational dynamics are adopted to be exactly Einsteinian with universal matter coupling. In the active sequence, \emph{adoption} means use of that established relativistic dynamics without claiming substrate derivation, while \emph{derivation} is reserved for a first-principles recovery from explicit substrate variables.

The present note stays strictly inside that boundary. It records only which kind of next manuscript would now count as an honest benchmark-facing gain at the current seed frontier.

\section{Fixed Inputs}

\begin{definition}[Seed-frontier fixed inputs]
The criterion recorded here uses the following fixed inputs.
\begin{enumerate}
    \item The benchmark exact-closure package, which fixes the public one-metric GR target of the repository.
    \item The benchmark gatekeeping layer, which already blocks same-output deeper-origin relays from counting as new front-line progress.
    \item The root-nucleus collapse theorem, which proves that the larger benchmark-facing root core collapses to the smaller minimal root exact-shell transport nucleus \(\RootNucleus\).
    \item The seed-dynamics theorem, which proves that the nucleus \(\RootNucleus\) itself is the exact-shell transport image of a more primitive nucleus-generating substrate seed dynamics package \(\SeedPackage\).
\end{enumerate}
\end{definition}

\begin{remark}
The present note does not replace those manuscripts. It records the routing consequence of reading them together under the active safeguard against same-output recursive depth drift.
\end{remark}

\section{Why the Seed-Dynamics Theorem Is the Frontier}

\begin{definition}[Benchmark-neutral seed restatement]
A \emph{benchmark-neutral seed restatement} is a manuscript that preserves the same benchmark one-metric output, the same ordinary matter route, the same settled observer-side closure verdict, and the same seed package \(\SeedPackage\) while merely relocating the unexplained layer deeper or rewriting the seed-level derivation without supplying a new compression, necessity, uniqueness, or comparably sharp routing gain.
\end{definition}

\begin{proposition}[The seed-dynamics theorem is the live frontier]
At the current stage of the primitive-reservoir observer-reduction line, the theorem deriving the minimal root exact-shell transport nucleus from nucleus-generating substrate seed dynamics remains the live benchmark-facing frontier.
\end{proposition}

\begin{proof}
That theorem changes benchmark-facing status at current scope because it removes the nucleus itself as the unexplained stopping point. The current benchmark-facing root datum is no longer merely a smaller internal root package; it is now shown to be the exact-shell transport image of a more primitive seed dynamics. By contrast, a manuscript that preserves the same benchmark package and the same seed package while merely pushing the unexplained layer deeper does not alter benchmark-facing status unless it adds a comparably sharp gain. Therefore the live frontier is now the seed-dynamics theorem rather than the earlier root-nucleus theorem or a later same-output relay.
\end{proof}

\begin{corollary}[Status of the former root-nucleus theorem]
The former root-nucleus collapse theorem remains preserved benchmark-facing main-line history, but under the current routing safeguard it is no longer by itself the default current frontier.
\end{corollary}

\begin{proof}
That theorem matters because it proved that the larger benchmark-facing root core contains no extra benchmark-relevant freedom beyond \(\RootNucleus\). But the later seed-dynamics theorem shows that even \(\RootNucleus\) is no longer the present unexplained stopping point. The former theorem is preserved history and handoff, while the seed-dynamics theorem defines the current frontier.
\end{proof}

\section{The Next Benchmark Criterion}

\begin{definition}[Admissible next benchmark-facing seed manuscript]
At the current seed frontier, a manuscript counts as an admissible next benchmark-facing move only if it does at least one of the following:
\begin{enumerate}
    \item derives or justifies the nucleus-generating substrate seed dynamics package \(\SeedPackage\) from still more primitive substrate structure in a way that does more than merely relocate the unexplained layer deeper;
    \item proves a genuinely smaller theorem-level collapse strictly inside the observer-relevant seed package \(\SeedPackage\);
    \item does either of the previous two while preserving the same benchmark one-metric package, the same ordinary matter route, and the same claim boundary.
\end{enumerate}
\end{definition}

\begin{theorem}[Next honest benchmark-facing task at seed scope]
The next honest benchmark-facing manuscript at the current seed frontier must either derive or justify the nucleus-generating substrate seed dynamics from still more primitive substrate structure in a benchmark-relevant way, or else prove a genuinely smaller collapse inside the observer-relevant seed package. A manuscript that only repeats the minimal root exact-shell transport nucleus, or only repeats the same seed package at a deeper level without such a new gain, does not satisfy the criterion.
\end{theorem}

\begin{proof}
By Proposition 1, the live frontier already lies at the seed package. Therefore a second manuscript below it must improve on that status rather than merely accompany it. A deeper-origin manuscript can do so only if it changes benchmark-facing status by more than depth alone. If it simply reproduces the same seed package through another relay, the routing rule classifies it as benchmark neutral. The remaining honest alternatives are therefore exactly those stated in the definition: either a more primitive derivation that changes benchmark-facing status, or a sharper internal collapse of the current observer-relevant seed package.
\end{proof}

\section{What the Next Move Should Not Be}

\begin{proposition}[Screened next moves]
At the current seed frontier, none of the following is the default next benchmark-facing move:
\begin{enumerate}
    \item another same-output restatement of the minimal root exact-shell transport nucleus;
    \item another same-output root, foundation, or ground relay;
    \item another same-output restatement of the same nucleus-generating substrate seed dynamics package;
    \item a reopening of stationary-reduction, stationary-Hessian, spectral, or selector layers;
    \item a manuscript that introduces a second operative metric, enlarges the ordinary matter law, or uses stronger promotion language.
\end{enumerate}
\end{proposition}

\begin{proof}
Items (1)--(3) fail the preceding criterion because they do not directly address the current seed burden. They revisit either an already-compressed older presentation or the same seed-level package rather than removing the remaining seed-level freedom or proving a smaller internal core. Item (4) likewise does not directly address the live burden at current scope. It reopens lower or screened layers as if they were still frontier-defining rather than settling the present seed-level task. Item (5) violates the fixed claim boundary of the benchmark package and therefore cannot count as a disciplined next step on the active line. The benchmark package remains the exact one-metric GR package already fixed by adoption, so any admissible next manuscript must preserve that boundary.
\end{proof}

\begin{remark}
This screening statement is not a denial that those lower-layer manuscripts exist or matter historically. It is a routing statement: they are not the default way to spend the next benchmark-facing move from the current seed frontier.
\end{remark}

\section{Conclusion}

The positive result of this note is routing discipline at the current seed frontier. The new seed-dynamics theorem remains the live benchmark-facing gain because it removes the minimal root exact-shell transport nucleus as the unexplained stopping point on the active line. The former root-nucleus collapse theorem remains preserved main-line history, but under the recursive-depth safeguard it is no longer itself the default current frontier.

The scope remains narrow and honest. The benchmark package is unchanged: the one operative metric is still exactly \(\CompletedMetric\), the ordinary matter route is unchanged, there is no second operative metric, and the low-energy matter law is not enlarged. Nothing here upgrades adoption to derivation or licenses stronger promotion language.

The next open burden is therefore clear. The next benchmark-facing manuscript should derive or justify the nucleus-generating substrate seed dynamics from still more primitive substrate structure in a way that changes benchmark-facing status, unless the sharper honest gain available instead is a genuinely smaller theorem-level collapse inside the observer-relevant seed package itself. That is the clean next manuscript target at current scope.

\end{document}
